\documentclass{peri}
\usepackage[utf8]{inputenc}
\usepackage[T1]{fontenc}
\usepackage{geometry}
\usepackage{graphicx}
\usepackage{float}
\usepackage{hyperref}
\usepackage[english,indonesian]{babel}

\title{Optimasi Pemrosesan Template LaTeX Menggunakan Pendekatan Berbasis AI}
\author{Academic Support Agent}
\date{\today}

\begin{document}
\maketitle

\begin{abstract}

\end{abstract}

\chapter{PENDAHULUAN}


\chapter{TINJAUAN PUSTAKA}


% Original content (commented out):
% % Template for PERI - Revista de Filosofia
% ISSN 2175-1811, Florianópolis, Santa Catarina, Brasil
%e-mail: revistaperi@contato.ufsc.br 
%site: http://www.nexos.ufsc.br/index.php/peri

% This template contains the following files: 
% peri.cls - peri's document class (please don't modify things there);
% template.tex - type your paper here;
% template.pdf - things should look like this;
% references - bibliography file 

\documentclass{peri}

\usepackage[english]{babel} % define the language of your article. Compatible with: portuguese, english, spanish and french.

%required packages:
%\usepackage[utf8]{inputenc} %if you are not using XeLaTeX as a compiler, enable this package.
\usepackage[T1]{fontenc}
\usepackage{amssymb}
\usepackage{amsthm}
\usepackage[alf,abnt-emphasize=bf]{abntex2cite}
\usepackage{lipsum}
%custom packages: please, keep it simple

%if needed, define your custom ambient for axioms, theorems, etc.
%\newtheorem{axiom}{Axiom}
%\newtheorem{theorem}{Theorem}
%\newtheorem{definition}{Definition}

%title of your paper
\title{Writting for PERI in \LaTeX}

% Author information (please, leave blank)
\author{Please leave blank\thanks{Affiliation; Funding info. E-mail.}}

\begin{document}

\maketitle

\begin{resumo}
Abstract in the main language of the article. 
\end{resumo}

\begin{palavraschave} Word1, Word 2, Word 3
\end{palavraschave}

\begin{abstract}
Abstract in portuguese. (If the main language is portuguese, write this abstract in english. \end{abstract}

\begin{keywords}
Palavra1, Palavra2, Palavra3.
\end{keywords}

\section*{First section}

Sections are not numbered. There are no subsections or subsubsections allowed. Please use the ``*'' character to organize your text \begin{verbatim}
    \section*{}
\end{verbatim}

References follow the \texttt{abntex2cite} package.\footnote{This is a footnote. See https://www.ctan.org/pkg/abntex2 for more information on the package.} For short quotations (up to 3 lines): ``text up to three lines (in the .pdf file, not here)'' \cite[p.~1]{template_test}. If you don't want to mention the page: \cite{template_test}. To compile the author's surname, the year and/or page of the publication, use the command: \citeonline[p.~32]{template_test}. For long quotations, use the \texttt{quotation} ambient.

\begin{quotation}
\lipsum[3]\cite[pp.~1--2]{template_test}.
\end{quotation}

\section*{Second section}

This is the second section. \textit{This is an emphasis}. 

\subsection*{Subsection}
This is a subsection.
\subsubsection*{Subsubsection}
This is a subsubsection.

\newpage
\bibliography{references} %insert the name of your .bib file here

\end{document}
% well done!

\end{document}